% Generated draft: requires mathematical and source review.
\documentclass[11pt]{article}
\usepackage[margin=25mm]{geometry}
\usepackage{fontspec}
\setmainfont{Latin Modern Roman}
\usepackage{amsmath,amssymb,mathtools,unicode-math}
\setmathfont{Latin Modern Math}
\usepackage{enumitem,needspace,xurl,hyperref}
\setlength{\parindent}{0pt}
\setlength{\parskip}{5pt}
\setlength{\emergencystretch}{3em}
\newcommand{\N}{\mathbb{N}}
\newcommand{\Z}{\mathbb{Z}}
\newcommand{\Q}{\mathbb{Q}}
\newcommand{\R}{\mathbb{R}}
\newcommand{\C}{\mathbb{C}}
\newcommand{\F}{\mathbb{F}}
\newcommand{\A}{\mathbb{A}}
\newcommand{\PP}{\mathbb{P}}
\newcommand{\E}{\mathbb{E}}
\newcommand{\eps}{\varepsilon}
\newcommand{\dd}{\,\mathrm{d}}
\begin{document}
\section*{UnsolvedMath ID: 30000817}
\textbf{Source problem code:} catalogue entry
\subsection{Definitions and mathematical statement}
Let \(P\) be a lattice polytope of dimension \(d\) in a lattice \(M\). Work relative to its affine lattice \(\operatorname{aff}(P)\cap M\), with difference lattice \(M_P=\operatorname{span}_{\R}(P-P)\cap M\). A full-dimensional simplex \(\Delta=\operatorname{conv}(v_0,\ldots,v_d)\subseteq P\), with lattice vertices, is unimodular if \(v_1-v_0,\ldots,v_d-v_0\) form a \(\Z\)-basis of \(M_P\). A unimodular covering is a finite collection of such simplices satisfying
\[
 P=\bigcup_{i=1}^N\Delta_i.
\]
No face-to-face condition on intersections is imposed for a cover. A unimodular triangulation is a triangulation by unimodular simplices; regularity additionally means that the triangulation is induced by a lifting height function. In particular, every regular unimodular triangulation of \(P\) gives a unimodular covering of \(P\).

Under the usual integer-decomposition convention, normality means
\[
 \forall m\geq1\ \forall x\in mP\cap M,\qquad
 x=x_1+\cdots+x_m\text{ for some }x_i\in P\cap M.
\]
The supplied record names two conditions without defining them: \(A(P)\), ``pairwise face unimodularity,'' reported to imply normality, and \(B(P)\), ``facet unimodularity,'' reported to imply a regular unimodular triangulation. It asks the two implications
\[
 A(P)\ \Longrightarrow\ P\text{ has a unimodular covering}\ ?
 \qquad
 B(P)\ \Longrightarrow\ P\text{ has a unimodular covering}\ ?
\]
Here \(A\) and \(B\) are labels for missing source-specific hypotheses, not definitions recovered from their names.

As literally written, if the stated consequence of \(B(P)\) is a regular unimodular triangulation of the whole polytope \(P\), its covering consequence is immediate from the definitions. Therefore the second purported open implication cannot be a separate unresolved implication under that reading. Possible differences involving triangulations of facets rather than \(P\), or another meaning of the hypothesis, require source review; none is selected here. Normality alone is not substituted for the undefined condition \(A(P)\). The workshop title and dated assessment mention smooth polytopes, but the displayed extraction does not explicitly require smoothness.
\subsection{Short English statement}
The record asks whether two face conditions yield a cover by lattice simplices of normalized volume one. The conditions are undefined in the extraction, and the second question is automatic if its stated premise already provides a unimodular triangulation of the whole polytope. This discrepancy prevents a reliable complete reconstruction.
\subsection{Sources}
\begin{itemize}
\item Local UnsolvedMath registry, permanent ID \(30000817\), exact source code \texttt{catalogue entry}, proposed year 2007; no individual proposer is supplied.
\item \emph{Mini-Workshop: Projective Normality of Smooth Toric Varieties (2008)}, Oberwolfach Reports, DOI 10.4171/owr/2007/39, pp. 2283--2320. \url{https://doi.org/10.4171/owr/2007/39}.
\item Supplied source labelled ``covering-properties thesis'': \url{https://refubium.fu-berlin.de/bitstream/handle/fub188/26773/Thesis_Codenotti.pdf?sequence=3}.
\item The supplied \texttt{corrected\_verified} note says both sufficient conditions and their shared covering conclusion are explicit in the supplied statement. The assessment checked 2026-08-21 labels the problem open and describes unresolved covering questions and the general smooth-polytope cover problem. This is dated supplied triage; it does not resolve the logical and definitional issues above.
\end{itemize}
\end{document}
